\chapter{Espacios vectoriales normados}\label{ch:b2:nvs}

Cuando el \omterm{def:b2:metric:def}{espacio métrico} es un espacio vectorial y la distancia
procede de una \omterm{def:b2:nvs:norm}{norma}, la topología y el álgebra lineal empiezan a
interactuar: las aplicaciones lineales son \omterm{def:b2:metric:continuity}{continuas} exactamente
cuando están acotadas sobre la bola unidad, la dimensión finita obliga
a todas las \omterm{def:b2:nvs:norm}{normas} a coincidir y la \omterm{def:b2:metric:complete}{completitud} convierte las series
absolutamente convergentes en convergentes. La divisoria entre
dimensión finita e infinita ---cristalizada en el teorema de Riesz---
es la lección más profunda del capítulo.

En todo el capítulo, $E, F$ son espacios vectoriales sobre $K = \R$ o
$\C$.

\section{Normas}

\begin{definition}\label{def:b2:nvs:norm}
Una \emph{norma}\index{norma} sobre $E$ es una aplicación
$\norm{\,\cdot\,} \colon E \to \R_+$ tal que, para todos $x, y \in
E$ y $\lambda \in K$:
\[
\norm x = 0 \iff x = 0,
\qquad
\norm{\lambda x} = \abs\lambda\,\norm x,
\qquad
\norm{x + y} \leq \norm x + \norm y .
\]
Entonces $d(x, y) = \norm{x - y}$ es una distancia y se aplica todo
el~\cref{ch:b2:metric}. La desigualdad triangular inversa
$\bigl|\norm x - \norm y\bigr| \leq \norm{x - y}$ hace que la propia
norma sea $1$-lipschitziana; la suma y la multiplicación por
escalares son \omterm{def:b2:metric:continuity}{continuas} (estimaciones $\norm{(x + y) - (x' + y')}
\leq \norm{x - x'} + \norm{y - y'}$, etc.).
\end{definition}

\begin{example}\label{ex:b2:nvs:examples}
Sobre $K^n$:
\[
\norm{x}_1 = \sum_i \abs{x_i},
\qquad
\norm{x}_2 = \Bigl(\sum_i \abs{x_i}^2\Bigr)^{1/2},
\qquad
\norm{x}_\infty = \max_i \abs{x_i}
\]
($\norm\cdot_2$ es una \omterm{def:b2:nvs:norm}{norma} por Cauchy--Schwarz, volumen del primer
año). Sobre $C(\intcc{a}{b})$:
\[
\norm f_\infty = \sup \abs f,
\qquad
\norm f_1 = \int_a^b \abs f,
\qquad
\norm f_2 = \Bigl(\int_a^b \abs f^2\Bigr)^{1/2},
\]
siendo las dos últimas \omterm{def:b2:nvs:norm}{normas} gracias a la positividad estricta de
la integral y a la desigualdad integral de Cauchy--Schwarz (volumen
del primer año). Sobre matrices: cualquier \omterm{def:b2:nvs:norm}{norma} sobre
$\mathcal M_n(K) \simeq K^{n^2}$; las \omterm{thm:b2:nvs:continuouslinear}{normas de operador} de más
abajo son las estructuralmente importantes.
\end{example}

\begin{definition}[Normas equivalentes]\label{def:b2:nvs:equivalent}
Dos \omterm{def:b2:nvs:norm}{normas} $N_1, N_2$ sobre $E$ son
\emph{equivalentes}\index{normas equivalentes} cuando existen
constantes $c, C > 0$ con
\[
c\,N_1 \leq N_2 \leq C\, N_1 .
\]
Dos \omterm{def:b2:nvs:norm}{normas} equivalentes tienen los mismos \omterm{def:b2:metric:topology}{abiertos}, las mismas
sucesiones convergentes y de Cauchy, y los mismos subconjuntos
\omterm{def:b2:metric:compact}{compactos} y \omterm{def:b2:metric:complete}{completos}: el mismo análisis.
\end{definition}

\begin{example}[No equivalencia en dimensión infinita]\label{ex:b2:nvs:nonequivalent}
Sobre $C(\intcc{0}{1})$ se cumple siempre $\norm f_1 \leq \norm
f_\infty$, pero no hay ninguna cota recíproca: $f_n(x) = x^n$ tiene
$\norm{f_n}_\infty = 1$ y $\norm{f_n}_1 = \frac{1}{n+1} \to 0$. Así
pues, $f_n \to 0$ para $\norm\cdot_1$ pero no para
$\norm\cdot_\infty$: las dos \omterm{def:b2:nvs:norm}{normas} discrepan sobre la propia noción
de convergencia.
\end{example}

\begin{example}[Constantes explícitas en dimensión $n$]\label{ex:b2:nvs:constants}
Sobre $K^n$, las tres \omterm{def:b2:nvs:norm}{normas} clásicas son \omterm{def:b2:nvs:equivalent}{equivalentes} con
constantes óptimas:
\[
\norm x_\infty \leq \norm x_2 \leq \norm x_1
\leq \sqrt n\,\norm x_2 \leq n\,\norm x_\infty ,
\]
donde la cota central $\norm x_1 \leq \sqrt n\norm x_2$ procede de
Cauchy--Schwarz contra el vector de unos. Vectores extremales:
$e_1$ convierte las dos primeras desigualdades en igualdades, y
$(1, 1, \dots, 1)$ las dos últimas. La dimensión $n$ aparece
visiblemente en las constantes: es la semilla cuantitativa del
fracaso en dimensión infinita, pues cuando $n \to \infty$ no
sobrevive ninguna constante uniforme, que es exactamente lo que
exhibe el~\cref{ex:b2:nvs:nonequivalent} sobre espacios de
funciones.
\end{example}

\begin{omfigure}
\begin{tikzpicture}[scale=1.5]
  \draw[->] (-1.45,0) -- (1.5,0) node[below, font=\small] {$x$};
  \draw[->] (0,-1.45) -- (0,1.5) node[left, font=\small] {$y$};
  \draw[omProp, very thick]
    (1,1) -- (-1,1) -- (-1,-1) -- (1,-1) -- cycle;
  \draw[omThm, very thick] (0,0) circle (1);
  \draw[omDef, very thick]
    (1,0) -- (0,1) -- (-1,0) -- (0,-1) -- cycle;
  \node[omDef, font=\small] at (0.36,0.28)
    {$\norm\cdot_1$};
  \node[omThm, font=\small] at (0.55,0.92)
    {$\norm\cdot_2$};
  \node[omProp, font=\small] at (1.24,1.13)
    {$\norm\cdot_\infty$};
\end{tikzpicture}

{\small Las bolas unidad de las tres \omterm{def:b2:nvs:norm}{normas} clásicas de $\R^2$,
encajadas tal como dictan las desigualdades del~\cref{ex:b2:nvs:constants}: bola más pequeña, \omterm{def:b2:nvs:norm}{norma} más grande. La
redondez importa: los lados planos del rombo y del cuadrado son
exactamente los fallos de convexidad estricta que se explotan en el
problema de fin de semana del~\cref{ch:b2:realfun} y en el de este
capítulo (pregunta 4).}
\end{omfigure}

\section{Aplicaciones lineales continuas}

\begin{theorem}[Caracterización]\label{thm:b2:nvs:continuouslinear}
Para una aplicación lineal $u \colon E \to F$ entre espacios
normados, las afirmaciones siguientes son \omterm{def:b2:nvs:equivalent}{equivalentes}:
\begin{enumerate}
  \item $u$ es \omterm{def:b2:metric:continuity}{continua};
  \item $u$ es \omterm{def:b2:metric:continuity}{continua} en $0$;
  \item $u$ está acotada sobre la bola unidad cerrada: $\sup_{\norm
        x \leq 1} \norm{u(x)} < \infty$;
  \item existe $C \geq 0$ con $\norm{u(x)} \leq C \norm x$ para todo
        $x$;
  \item $u$ es \omterm{def:b2:metric:continuity}{lipschitziana}.
\end{enumerate}
La menor de esas constantes $C$ es la \emph{norma de
operador}\index{norma de operador} $\vertiii{u} = \sup_{\norm x \leq
1}\norm{u(x)} = \sup_{x \neq 0} \frac{\norm{u(x)}}{\norm x}$; hace
del espacio $\mathcal{L}_c(E, F)$ de las aplicaciones lineales
\omterm{def:b2:metric:continuity}{continuas} un espacio normado, con
\[
\vertiii{v \circ u} \leq \vertiii v\, \vertiii u .
\]
\end{theorem}

\begin{proof}
(1 $\Rightarrow$ 2) es trivial. (2 $\Rightarrow$ 3): la \omterm{def:b2:metric:continuity}{continuidad}
en $0$ con $\varepsilon = 1$ da un $\delta$ con $\norm x \leq \delta
\Rightarrow \norm{u(x)} \leq 1$; la homogeneidad reduce entonces
cualquier $x$ con $\norm x \leq 1$ dentro de esa bola y lo devuelve:
\[
\norm{u(x)} = \frac1\delta\,\norm{u(\delta x)} \leq
\frac1\delta ,
\]
ya que $\norm{\delta x} \leq \delta$. (3 $\Rightarrow$ 4): para $x
\neq 0$, aplíquese la cota a $\frac{x}{\norm x}$. (4 $\Rightarrow$
5): $\norm{u(x) - u(y)} = \norm{u(x - y)} \leq C\norm{x - y}$. (5
$\Rightarrow$ 1) es conocido.

Axiomas de \omterm{def:b2:nvs:norm}{norma} para $\vertiii\cdot$: la homogeneidad y la
separación son claras ($\vertiii u = 0$ obliga a $u = 0$ sobre la
bola y, por tanto, en todas partes); la desigualdad triangular sale
de $\norm{(u + v)(x)} \leq \norm{u(x)} + \norm{v(x)}$.
Submultiplicatividad: $\norm{v(u(x))} \leq \vertiii v\,\norm{u(x)}
\leq \vertiii v \vertiii u \norm x$.
\end{proof}

\begin{example}\label{ex:b2:nvs:operatornorms}
Sobre $\bigl(C(\intcc{0}{1}), \norm\cdot_\infty\bigr)$: la
evaluación $f \mapsto f(0)$ tiene \omterm{thm:b2:nvs:continuouslinear}{norma de operador} $1$; la
integración $f \mapsto \int_0^1 f$ tiene \omterm{def:b2:nvs:norm}{norma} $1$; la aplicación
$f \mapsto \int_0^1 t f(t)\dd t$ tiene \omterm{def:b2:nvs:norm}{norma} $\int_0^1 t\,\dd t =
\frac12$ (cota superior por la desigualdad triangular para
integrales, alcanzada en $f \equiv 1$). Pero la derivación, de
$(C^1, \norm\cdot_\infty)$ en $(C^0, \norm\cdot_\infty)$, \emph{no}
es \omterm{def:b2:metric:continuity}{continua}: $\norm{\sin(nx)}_\infty = 1$ mientras que la derivada
tiene \omterm{def:b2:nvs:norm}{norma} del supremo $n$. Lineal no implica \omterm{def:b2:metric:continuity}{continua} en
dimensión infinita.
\end{example}

\begin{example}[Dos normas, dos veredictos sobre una misma sucesión]\label{ex:b2:nvs:sqrtnxn}
Sobre $C(\intcc01)$, sea $g_n(x) = \sqrt{n}\,x^n$. Entonces
\[
\norm{g_n}_1 = \frac{\sqrt n}{n + 1} \longrightarrow 0,
\qquad
\norm{g_n}_2^2 = \frac{n}{2n + 1} \longrightarrow \frac12,
\qquad
\norm{g_n}_\infty = \sqrt n \longrightarrow \infty :
\]
una sucesión, tres \omterm{def:b2:nvs:norm}{normas}, tres comportamientos ---convergencia a
cero, ausencia de convergencia (las \omterm{def:b2:nvs:norm}{normas} se estabilizan en
$\frac1{\sqrt2}$ mientras que el límite puntual es $0$) y
explosión---. La masa que se concentra cerca de $x = 1$ es
invisible para $\norm\cdot_1$, medio visible para $\norm\cdot_2$ y
dominante para $\norm\cdot_\infty$. En dimensión infinita, ``¿es
convergente?'' no es una pregunta sobre una sucesión: es una
pregunta sobre una sucesión \emph{y} una \omterm{def:b2:nvs:norm}{norma}.
\end{example}

\begin{method}[Cómo calcular una norma de operador]\label{met:b2:nvs:opnorm}
Siempre en dos movimientos. \emph{Cota superior:} estímese
$\norm{u(x)}$ por $C\norm x$ mediante desigualdades triangulares,
Cauchy--Schwarz o cotas integrales; eso demuestra $\vertiii u \leq
C$. \emph{Testigo:} exhíbase o bien un $x_0 \neq 0$ concreto con
$\norm{u(x_0)} = C\norm{x_0}$ (la cota se alcanza), o bien una
sucesión de vectores unitarios $x_n$ con $\norm{u(x_n)} \to C$ (la
cota se aproxima). Ambos movimientos son obligatorios: una cota
superior sola da únicamente $\vertiii u \leq C$, y un testigo solo,
únicamente $\vertiii u \geq C$. En dimensión infinita el testigo
puede tener que ser una sucesión, pues el supremo no tiene por qué
alcanzarse (\cref{exo:b2:nvs:8}).
\end{method}

\begin{example}[Los operadores diagonales ven todas las normas igual]\label{ex:b2:nvs:diagonalnorm}
Para $D = \operatorname{diag}(d_1, \dots, d_n)$ sobre $K^n$ con
\emph{cualquiera} de las \omterm{def:b2:nvs:norm}{normas} $\norm\cdot_1, \norm\cdot_2,
\norm\cdot_\infty$: de $\abs{d_ix_i} \leq
\bigl(\max_j\abs{d_j}\bigr)\abs{x_i}$ coordenada a coordenada se
sigue $\norm{Dx} \leq \max_j\abs{d_j}\,\norm x$; y $x = e_{j_0}$
(un índice maximizante) la alcanza. Luego $\vertiii D =
\max_j\abs{d_j}$ en los tres casos: para las aplicaciones
diagonales, todas las \omterm{def:b2:nvs:norm}{normas} razonables cuentan la misma
historia, el mayor factor de dilatación. Todo lo difícil de las
\omterm{thm:b2:nvs:continuouslinear}{normas de operador} tiene que ver con el comportamiento \emph{no}
diagonal, y por eso las \omterm{def:b2:nvs:norm}{normas} adaptadas del problema de fin de
semana de este capítulo (pregunta 22) funcionan obligando antes a
la matriz a hacerse diagonal.
\end{example}

\begin{example}[Sumas por columnas: el gemelo en norma $1$ del~\cref{exo:b2:nvs:4}]\label{ex:b2:nvs:columnsums}
Sobre $(\R^n, \norm\cdot_1)$, la \omterm{thm:b2:nvs:continuouslinear}{norma de operador} de una matriz
$A$ es la mayor suma de valores absolutos por \emph{columnas}.
Apliquemos el método: para $\norm x_1 \leq 1$,
\[
\norm{Ax}_1 = \sum_i\Bigl|\sum_j a_{ij}x_j\Bigr|
\leq \sum_j \abs{x_j}\sum_i\abs{a_{ij}}
\leq \Bigl(\max_j\sum_i\abs{a_{ij}}\Bigr)\norm x_1 ,
\]
y la cota se alcanza en $x = e_{j_0}$ para una columna maximizante
$j_0$: el testigo más pulcro imaginable. Así, para $A =
\begin{psmallmatrix}1 & -2\\ 3 & 1\end{psmallmatrix}$ resulta
$\vertiii A_1 = \max(1 + 3,\ 2 + 1) = 4$, mientras que también
$\vertiii A_\infty = 4$ (por filas), coincidencia en este caso y no
ley: altérense asimétricamente las entradas y las dos \omterm{def:b2:nvs:norm}{normas} se
separan. Filas para $\norm\cdot_\infty$, columnas para
$\norm\cdot_1$: la regla mnemotécnica es que los vectores unitarios
de cada \omterm{def:b2:nvs:norm}{norma} (patrones de signos o vectores de la base,
respectivamente) seleccionan las sumas correspondientes.
\end{example}

\begin{proposition}[Aplicaciones bilineales]\label{prop:b2:nvs:bilinear}
Una aplicación bilineal $b \colon E \times F \to G$ es \omterm{def:b2:metric:continuity}{continua} si
y solo si $\norm{b(x,y)} \leq C\norm x\,\norm y$ para cierta $C$; y
entonces es \omterm{def:b2:metric:continuity}{lipschitziana} sobre los acotados. (Mismo esquema de
demostración; el producto $(u, v) \mapsto v \circ u$ y la
multiplicación de matrices son los ejemplos clave.)
\end{proposition}

\begin{proof}
Si se cumple la cota:
\[
b(x,y) - b(x_0,y_0) = b(x - x_0,\, y) + b(x_0,\, y - y_0),
\]
luego $\norm{b(x,y) - b(x_0,y_0)} \leq C\norm{x - x_0}\norm y +
C\norm{x_0}\norm{y - y_0}$: hay \omterm{def:b2:metric:continuity}{continuidad} en $(x_0, y_0)$ y una
cota \omterm{def:b2:metric:continuity}{lipschitziana} donde $\norm x, \norm y \leq R$. Recíprocamente,
la \omterm{def:b2:metric:continuity}{continuidad} en $(0,0)$ da un $\delta$ con $\norm{b(x,y)} \leq 1$
para $\norm x, \norm y \leq \delta$; reescálense ambas variables.
\end{proof}

\section{Dimensión finita}

\begin{theorem}[Equivalencia de normas en dimensión finita]\label{thm:b2:nvs:finitedim}
Sobre un espacio de dimensión finita, \emph{todas las \omterm{def:b2:nvs:norm}{normas} son
\omterm{def:b2:nvs:equivalent}{equivalentes}}.\index{equivalencia de normas} En consecuencia, en
dimensión finita la convergencia, el ser \omterm{def:b2:metric:topology}{abierto}, la \omterm{def:b2:metric:compact}{compacidad} y la
\omterm{def:b2:metric:complete}{completitud} son nociones independientes de la \omterm{def:b2:nvs:norm}{norma}; \omterm{def:b2:metric:compact}{compacto} $=$
cerrado y acotado; el espacio es \omterm{def:b2:metric:complete}{completo}; y toda aplicación lineal
(o multilineal) \emph{desde} un espacio de dimensión finita es
\omterm{def:b2:metric:continuity}{continua}.
\end{theorem}

\begin{proof}
Fijemos una base e identifiquemos $E \simeq K^n$; basta comparar una
\omterm{def:b2:nvs:norm}{norma} cualquiera $N$ con $\norm\cdot_\infty$.

\emph{Una dirección es álgebra:} $N(x) = N(\sum x_i e_i) \leq \sum
\abs{x_i} N(e_i) \leq C \norm x_\infty$ con $C = \sum N(e_i)$. Esto
muestra además que $N$ es \omterm{def:b2:metric:continuity}{continua} sobre $(K^n,
\norm\cdot_\infty)$ (es $C$-lipschitziana: $\abs{N(x) - N(y)} \leq
N(x - y)$).

\emph{La otra es topología:} la esfera unidad $S = \{x :
\norm{x}_\infty = 1\}$ es cerrada y acotada en $(K^n,
\norm\cdot_\infty)$ y, por tanto, \omterm{def:b2:metric:compact}{compacta}
(\cref{thm:b2:metric:compactprops}~(2), válido para $\C^n \simeq
\R^{2n}$). La función \omterm{def:b2:metric:continuity}{continua} $N$ alcanza su mínimo $c$ sobre $S$;
y $c > 0$, ya que $N$ solo se anula en $0 \notin S$. La homogeneidad
extiende la cota: $N(x) \geq c \norm{x}_\infty$ para todo $x$.

Consecuencias: todos los enunciados se reducen a $(K^n,
\norm\cdot_\infty)$, donde son conocidos
(\cref{thm:b2:metric:rncomplete},
\cref{thm:b2:metric:compactprops}); y una aplicación lineal $u$
desde un $E$ de dimensión finita cumple $\norm{u(x)} \leq
\sum\abs{x_i} \norm{u(e_i)} \leq C'\norm{x}_\infty$: la cota (4) del~\cref{thm:b2:nvs:continuouslinear}.
\end{proof}

\begin{corollary}\label{cor:b2:nvs:closedsubspace}
Todo subespacio de dimensión finita de un espacio normado es
cerrado.
\end{corollary}

\begin{proof}
Es \omterm{def:b2:metric:complete}{completo} para la \omterm{def:b2:nvs:norm}{norma} inducida
(\cref{thm:b2:nvs:finitedim}), y los subconjuntos \omterm{def:b2:metric:complete}{completos} son
cerrados (\cref{def:b2:metric:complete}).
\end{proof}

\begin{example}[Una mejor aproximación calculada por simetría]\label{ex:b2:nvs:bestaffine}
En $\bigl(C(\intcc{-1}{1}), \norm\cdot_\infty\bigr)$, ¿a qué
distancia está $f(x) = \abs x$ del subespacio (cerrado, de
dimensión dos) de las funciones afines $a + bx$? Por simetría,
sustituir $a + bx$ por $a - bx$ no altera $\norm{f - (a \pm
bx)}_\infty$, y el punto medio $a$ lo hace al menos igual de bien
(desigualdad triangular sobre la media): basta considerar las
constantes. Para una constante $a$:
\[
\norm{\abs x - a}_\infty = \max\,(1 - a,\ a)
\geq \frac12 ,
\]
mínimo en $a = \frac12$: la distancia es $\frac12$, alcanzada por
la constante $\frac12$. Obsérvese la curva de error $\abs x -
\frac12$: alcanza $\pm\frac12$ alternativamente en $x = -1, 0, 1$,
es decir, tres extremos de signo alterno para una mejor
aproximación desde una familia de dos parámetros. Ese patrón de
\emph{equioscilación} no es casual: es la firma de la optimalidad
que el problema de fin de semana de este capítulo convierte en el
teorema de Chebyshev.
\end{example}

\begin{example}[Subespacios cerrados frente a densos]\label{ex:b2:nvs:closeddense}
En $E = \bigl(C(\intcc{0}{1}), \norm\cdot_\infty\bigr)$, cada
$\R_n[X]$ (polinomios de grado $\leq n$, restringidos a
$\intcc01$) es un subespacio de dimensión finita y, por tanto,
\emph{cerrado}: un límite uniforme de polinomios de grado $\leq n$
lo es también. Pero la unión $\R[X]$ de todos ellos es \emph{densa}
en $E$ (teorema de aproximación de Weierstrass, demostrado en el~\cref{ch:b2:funcseq}), y los \omterm{def:b2:reduction:eigen}{subespacios propios} densos son todo lo
contrario de cerrados. Moraleja: el ser cerrado es un privilegio de
los subespacios de dimensión finita; apilando pisos cerrados se
puede construir un rascacielos denso.
\end{example}

\begin{theorem}[Riesz]\label{thm:b2:nvs:riesz}
La bola unidad cerrada de un espacio normado $E$ es \omterm{def:b2:metric:compact}{compacta}
\emph{si y solo si} $\dim E < \infty$.\index{teorema de Riesz}
\end{theorem}

\begin{proof}
Dimensión finita: basta con ser cerrada y acotada
(\cref{thm:b2:nvs:finitedim}).

Recíprocamente, supongamos $\dim E = \infty$. \emph{Lema de Riesz:}
para todo subespacio cerrado propio $F \subsetneq E$ y todo
$\varepsilon \in \intoo{0}{1}$ existe un vector unitario $x$ con
$d(x, F) \geq 1 - \varepsilon$. Demostración: tómese $y \notin F$,
sea $\delta = d(y, F) > 0$ ($F$ es cerrado), elíjase $f \in F$ con
$\norm{y - f} \leq \frac{\delta}{1 - \varepsilon}$ y póngase $x =
\frac{y - f}{\norm{y - f}}$: para todo $g \in F$,
\[
\norm{x - g} = \frac{\norm{y - (f + \norm{y-f}\,g)}}{\norm{y - f}}
\geq \frac{\delta}{\norm{y-f}} \geq 1 - \varepsilon ,
\]
pues el numerador es una distancia de $y$ a un punto de $F$.

Construyamos ahora vectores unitarios $x_1, x_2, \dots$ por
inducción: $F_k = \operatorname{Vect}(x_1, \dots, x_k)$ es de
dimensión finita y, por tanto, cerrado
(\cref{cor:b2:nvs:closedsubspace}) y propio; el lema de Riesz con
$\varepsilon = \frac12$ proporciona un unitario $x_{k+1}$ con
$d(x_{k+1}, F_k) \geq \frac12$. La sucesión cumple $\norm{x_p -
x_q} \geq \frac12$ para $p \neq q$: no tiene subsucesión
convergente, luego la bola unidad no es \omterm{def:b2:metric:compact}{compacta}.
\end{proof}

\begin{example}[Riesz como detector de dimensión]\label{ex:b2:nvs:rieszdetector}
¿Es $C(\intcc01)$ de dimensión finita? Riesz responde sin exhibir
ninguna familia libre infinita explícita: la sucesión $f_n(x) =
x^n$ está en la bola unidad cerrada y cumple, para $m > n$,
$\norm{f_n - f_m}_\infty \geq f_n(x_0) - f_m(x_0) > 0$ en puntos
adecuados, cosa que el problema de fin de semana de este capítulo
(pregunta 16) cuantifica con limpieza, obteniendo una subsucesión
cuyos términos están a distancia mutua $\geq \frac14$. No hay
subsucesión convergente, luego la bola no es \omterm{def:b2:metric:compact}{compacta} y $\dim
C(\intcc01) = \infty$ por el~\cref{thm:b2:nvs:riesz}. La
\omterm{def:b2:metric:compact}{compacidad} de la bola unidad es una dicotomía perfecta: se cumple
en dimensión finita y falla en dimensión infinita, sin terreno
intermedio; la geometría sola lee el tipo de dimensión.
\end{example}

\section{Espacios de Banach}

\begin{definition}\label{def:b2:nvs:banach}
Un \emph{espacio de Banach}\index{espacio de Banach} es un espacio
normado \omterm{def:b2:metric:complete}{completo}. Ejemplos: todo espacio normado de dimensión
finita (\cref{thm:b2:nvs:finitedim}); $\bigl(C(\intcc{a}{b}),
\norm\cdot_\infty\bigr)$ (\cref{thm:b2:metric:rncomplete});
$\mathcal{L}_c(E, F)$ si $F$ es de Banach (mismo esquema de
demostración que para las funciones \omterm{def:b2:metric:continuity}{continuas}). Contraejemplo:
$\bigl(C(\intcc{0}{1}), \norm\cdot_1\bigr)$ (\cref{exo:b2:nvs:7}).
\end{definition}

\begin{example}[La norma de operador de la integración]\label{ex:b2:nvs:integrationnorm}
Sobre $\bigl(C(\intcc01), \norm\cdot_\infty\bigr)$, sea $T(f)(x) =
\int_0^x f(t)\,\dd t$ (un endomorfismo: $T(f)$ es \omterm{def:b2:metric:continuity}{continua}).
Apliquemos el~\cref{met:b2:nvs:opnorm}. Cota superior:
\[
\abs{T(f)(x)} \leq \int_0^x\abs f \leq x\,\norm f_\infty \leq
\norm f_\infty ,
\]
luego $\vertiii T \leq 1$. Testigo: $f \equiv 1$ da $T(f)(x) = x$ y
$\norm{T(f)}_\infty = 1 = \norm f_\infty$: se alcanza y $\vertiii T
= 1$. Pero obsérvese que $\vertiii{T^2} = \frac12 \neq \vertiii
T^2$: en efecto, $T^2(f)(x) = \int_0^x(x - t)f(t)\dd t$ cumple
$\abs{T^2(f)(x)} \leq \frac{x^2}2\norm f_\infty$, de nuevo
alcanzado en $f \equiv 1$; y, en general, $\vertiii{T^n} =
\frac1{n!}$: la cota submultiplicativa $\vertiii T^n = 1$ se
equivoca por un factorial. Este es exactamente el fenómeno que el
truco de la iterada del problema de fin de semana del~\cref{ch:b2:metric} convierte en resolubilidad global de las
ecuaciones diferenciales lineales.
\end{example}

\begin{theorem}[Convergencia absoluta en espacios de Banach]\label{thm:b2:nvs:absoluteconvergence}
En un \omterm{def:b2:nvs:banach}{espacio de Banach}, si $\sum \norm{u_n} < \infty$ entonces
$\sum u_n$ converge y $\norm{\sum u_n} \leq \sum\norm{u_n}$. (La
teoría completa de las series en espacios normados es el~\cref{ch:b2:series}.)
\end{theorem}

\begin{proof}
Sumas parciales $S_N$: para $q > p$, $\norm{S_q - S_p} \leq
\sum_{n=p+1}^{q}\norm{u_n}$, que tiende a $0$ (criterio de Cauchy
para la serie real de las \omterm{def:b2:nvs:norm}{normas}); luego $(S_N)$ es de Cauchy y
por tanto converge. La desigualdad pasa al límite desde la
desigualdad triangular finita.
\end{proof}

\begin{example}[La exponencial de matrices, primer contacto]\label{ex:b2:nvs:matrixexp}
La \emph{exponencial de matrices}\index{exponencial de matrices}:
$\mathcal{M}_n(K)$ con cualquier \omterm{def:b2:nvs:norm}{norma} submultiplicativa
($\vertiii{AB} \leq \vertiii A \vertiii B$) es de Banach (dimensión
finita). Entonces, para toda $A$,
\[
\eu^A = \sum_{k=0}^{\infty} \frac{A^k}{k!}
\]
converge absolutamente ($\vertiii{A^k/k!} \leq \vertiii A^k /k!$,
sumable): está bien definida. El~\cref{ch:b2:diffeq} la explota
sistemáticamente.
\end{example}

\begin{example}[Una serie de Neumann que termina]\label{ex:b2:nvs:neumannfinite}
Para $A = \begin{psmallmatrix}0 & \frac12\\ 0 &
0\end{psmallmatrix}$: $\vertiii A < 1$ en cualquier
\omterm{thm:b2:nvs:continuouslinear}{norma de operador} construida sobre las \omterm{def:b2:nvs:norm}{normas} del~\cref{ex:b2:nvs:examples}, y $A^2 = 0$, de modo que la serie
geométrica se derrumba:
\[
(I - A)^{-1} = \sum_{k \geq 0} A^k = I + A =
\begin{pmatrix}1 & \tfrac12\\ 0 & 1\end{pmatrix},
\]
lo que se verifica con $(I - A)(I + A) = I - A^2 = I$. La
nilpotencia trunca la serie exactamente igual que truncaba la
exponencial en el~\cref{ch:b2:reduction}; y el ejemplo calibra las
expectativas: la inversa de Neumann es una serie infinita en
general, y un polinomio precisamente cuando la perturbación es
nilpotente, mientras que el error tras $N$ términos siempre está
acotado por la cola geométrica $\vertiii A^{N+1}/(1 - \vertiii A)$.
\end{example}

\begin{example}[La exponencial de un generador de rotaciones]\label{ex:b2:nvs:rotationexp}
Tomemos $A = \begin{psmallmatrix}0 & -\theta\\ \theta &
0\end{psmallmatrix}$. Entonces $A^2 = -\theta^2 I$, de modo que las
potencias cierran un \omterm{def:b2:structures:sn}{ciclo} de periodo cuatro y la serie se separa
en partes par e impar:
\[
\eu^{A} = \sum_{k}\frac{A^k}{k!}
= \Bigl(\sum_{j}\frac{(-1)^j\theta^{2j}}{(2j)!}\Bigr) I
+ \Bigl(\sum_{j}\frac{(-1)^j\theta^{2j+1}}{(2j+1)!}\Bigr)
\frac{A}{\theta}
= \begin{pmatrix}
\cos\theta & -\sin\theta\\
\sin\theta & \cos\theta
\end{pmatrix},
\]
donde todos los reagrupamientos están autorizados por la
convergencia absoluta. La exponencial de un generador antisimétrico
es una rotación, calculada aquí puramente a partir de la serie,
tres capítulos antes de que la ecuación diferencial $x' = Ax$
(\cref{ch:b2:diffeq}) explique \emph{por qué}: $\eu^{tA}$ es un
movimiento circular uniforme. La moraleja: las identidades entre
series de matrices se demuestran exactamente como las escalares,
una vez que una \omterm{def:b2:nvs:norm}{norma} submultiplicativa certifica la convergencia
absoluta.
\end{example}

\begin{example}[Norma del supremo quiere decir uniforme: el diccionario]\label{ex:b2:nvs:uniformdictionary}
El enunciado $\norm{f_n - f}_\infty \to 0$ \emph{es} la
convergencia uniforme: un solo número, $\sup_x\abs{f_n(x) -
f(x)}$, acota el error en todos los puntos a la vez. El
diccionario en acción sobre $f_n(x) = x^n$ en $\intcc{0}{1}$:
puntualmente, $f_n \to 0$ en $\intco{0}{1}$ y $f_n(1) = 1$; en
\omterm{def:b2:nvs:norm}{norma}, $\norm{f_n - 0}_\infty = 1 \not\to 0$, y en efecto el
límite puntual es discontinuo y por tanto inalcanzable como límite
en $\norm\cdot_\infty$ dentro de $C(\intcc01)$ (que es cerrado
para límites uniformes, \cref{thm:b2:metric:rncomplete}). Sobre
$\intcc{0}{a}$ con $a < 1$: $\norm{f_n}_\infty = a^n \to 0$; la
convergencia uniforme se restablece encogiendo el dominio. Todo
enunciado de convergencia del~\cref{ch:b2:funcseq} es un enunciado
sobre esta única \omterm{def:b2:nvs:norm}{norma}; tener el diccionario presente reduce a la
mitad aquel capítulo.
\end{example}

\begin{remark}[Errores frecuentes]\label{rem:b2:nvs:pitfalls}
(i) Una \omterm{thm:b2:nvs:continuouslinear}{norma de operador} depende de \emph{ambas} \omterm{def:b2:nvs:norm}{normas}
elegidas: una misma matriz tiene $\vertiii\cdot_\infty$ dada por
las sumas por filas (\cref{exo:b2:nvs:4}) y una
$\vertiii\cdot_1$ distinta (sumas por columnas); hablar de ``la''
\omterm{def:b2:nvs:norm}{norma} de una matriz sin nombrar las \omterm{def:b2:nvs:norm}{normas} subyacentes carece de
sentido. (ii) $\vertiii{AB} \leq \vertiii A\,\vertiii B$ es una
desigualdad, por lo general estricta: las potencias pueden
encogerse mucho más deprisa de lo que sugiere la cota $\vertiii
A^k$, que es todo el sentido de las \omterm{def:b2:nvs:norm}{normas} adaptadas (problema de
fin de semana de este capítulo, pregunta 22). (iii) ``Lineal
implica \omterm{def:b2:metric:continuity}{continua}'' es un privilegio de la dimensión finita: la
derivación sobre polinomios es lineal y no acotada
(\cref{ex:b2:nvs:operatornorms}). (iv) La convergencia absoluta de
$\sum u_n$ solo sirve de algo cuando el espacio es \omterm{def:b2:metric:complete}{completo} (el~\cref{exo:b2:series:9} construye el contraejemplo). (v) En
dimensión infinita, un supremo sobre la bola unidad es un supremo
de verdad: no hay que suponer que se alcanza
(\cref{exo:b2:nvs:8}).
\end{remark}

\begin{remark}[Perspectivas dentro de este volumen]\label{rem:b2:nvs:perspectives}
Quedan fijadas tres citas. Con el~\cref{ch:b2:series}: en un
\omterm{def:b2:nvs:banach}{espacio de Banach} las series absolutamente convergentes convergen,
de modo que las series geométrica y exponencial de operadores pasan
a ser herramientas cotidianas ---invertir $I - A$, definir
$\eu^{A}$ (\cref{ex:b2:series:neumann})---. Con el~\cref{ch:b2:funcseq} y el~\cref{ch:b2:powerseries}: la convergencia
de sucesiones de funciones y de series de potencias es la
convergencia en $\bigl(C, \norm\cdot_\infty\bigr)$
(\cref{ex:b2:nvs:uniformdictionary}), y el radio de convergencia es
un enunciado sobre qué series geométricas dominan. Con el~\cref{ch:b2:fourier}: las \omterm{def:b2:nvs:norm}{normas} $\norm\cdot_2$ y
$\norm\cdot_\infty$ discrepan de verdad sobre $C(\intcc{0}{1})$
(\cref{ex:b2:nvs:sqrtnxn}), y por eso la convergencia en media
cuadrática de las series de Fourier y la convergencia uniforme son
dos teoremas distintos con dos precios distintos.
\end{remark}

\begin{remark}[Dónde se usa este capítulo]
Las \omterm{thm:b2:nvs:continuouslinear}{normas de operador} y la serie geométrica mueven los argumentos
de perturbación del~\cref{ch:b2:diffcalc} (teorema de la función
inversa) y la \omterm{ex:b2:nvs:matrixexp}{exponencial de matrices} del~\cref{ch:b2:diffeq}; la
\omterm{thm:b2:nvs:finitedim}{equivalencia de normas} autoriza en silencio todo argumento del
tipo ``elige tu \omterm{def:b2:nvs:norm}{norma} favorita'' en el~\cref{ch:b2:funcseq} y más
allá; y la divisoria finito/infinito del teorema de Riesz ---hecha
cuantitativa en el problema de fin de semana de este capítulo--- es
la razón de que el volumen del tercer año necesite herramientas
nuevas (convergencia débil, Arzelà--Ascoli, proyecciones en
espacios de Hilbert) allí donde este volumen aún podía extraer
subsucesiones convergentes.
\end{remark}

\section{Ejercicios}

\begin{exercise}[$\star$]\label{exo:b2:nvs:1}
Sobre $\R^2$, dibuja las bolas unidad de $\norm\cdot_1$,
$\norm\cdot_2$ y $\norm\cdot_\infty$, y demuestra las desigualdades
$\norm x_\infty \leq \norm x_2 \leq \norm x_1 \leq 2\norm x_\infty$
con las mejores constantes en dimensión $2$.
\end{exercise}

\begin{exercise}[$\star$]\label{exo:b2:nvs:2}
¿Es $N(f) = \abs{f(0)} + \norm{f'}_\infty$ una \omterm{def:b2:nvs:norm}{norma} sobre
$C^1(\intcc{0}{1})$? Compárala con $\norm{f}_\infty$: una de las
desigualdades se cumple y la otra falla (exhíbelo).
\end{exercise}

\begin{exercise}[$\star$]\label{exo:b2:nvs:3}
Calcula la \omterm{thm:b2:nvs:continuouslinear}{norma de operador} de $u(f) = \int_0^1
f(t)\,\eu^t\,\dd t$ de $\bigl(C(\intcc{0}{1}),
\norm\cdot_\infty\bigr)$ en $\R$, y la del desplazamiento $S(x_1,
x_2, \dots, x_n) = (x_2, \dots, x_n, 0)$ sobre $(K^n,
\norm\cdot_\infty)$.
\end{exercise}

\begin{exercise}[$\star\star$]\label{exo:b2:nvs:4}
Sobre $(\R^n, \norm\cdot_\infty)$, demuestra que la
\omterm{thm:b2:nvs:continuouslinear}{norma de operador} de una matriz $A$ es $\vertiii A_\infty = \max_i
\sum_j \abs{a_{ij}}$ (la mayor suma de valores absolutos por filas).
Calcúlala para $\begin{pmatrix} 1 & -2\\ 3 & 1\end{pmatrix}$.
\end{exercise}

\begin{exercise}[$\star\star$]\label{exo:b2:nvs:5}
Demuestra que $GL_n(K)$ es \omterm{def:b2:metric:topology}{abierto} en $\mathcal{M}_n(K)$ y que $A
\mapsto A^{-1}$ es \omterm{def:b2:metric:continuity}{continua} sobre él. \emph{Indicación: para el
carácter \omterm{def:b2:metric:topology}{abierto}, si $\vertiii H < \frac{1}{\vertiii{A^{-1}}}$
entonces $A + H = A(I + A^{-1}H)$ con $\vertiii{A^{-1}H} < 1$, y $I
+ B$ es invertible para $\vertiii B < 1$ por la serie geométrica
(\cref{thm:b2:nvs:absoluteconvergence}); para la \omterm{def:b2:metric:continuity}{continuidad}, acota
$(A+H)^{-1} - A^{-1}$ con la misma serie.}
\end{exercise}

\begin{exercise}[$\star\star$]\label{exo:b2:nvs:6}
Sea $\varphi$ una forma lineal sobre un espacio normado $E$.
Demuestra que $\varphi$ es \omterm{def:b2:metric:continuity}{continua} si y solo si $\ker\varphi$ es
cerrado. \emph{(Si $\ker\varphi$ es cerrado y $\varphi \neq 0$, toma
$a$ con $\varphi(a) = 1$ y $r > 0$ con $B(a, r) \cap \ker\varphi =
\emptyset$; deduce $\abs{\varphi(h)} \leq \frac{1}{r}\norm h$
mediante un argumento de escala sobre $a -
\frac{h}{\varphi(h)}$.)}
\end{exercise}

\begin{exercise}[$\star\star$]\label{exo:b2:nvs:7}
Demuestra que $\bigl(C(\intcc{0}{1}), \norm\cdot_1\bigr)$ no es
\omterm{def:b2:metric:complete}{completo}: prueba que las funciones $f_n$, rampas afines de $0$ a
$1$ sobre $\bigl[\frac12 - \frac1n, \frac12\bigr]$ (con valor $0$
antes y $1$ después), forman una sucesión de Cauchy sin límite
\omterm{def:b2:metric:continuity}{continuo} para $\norm\cdot_1$.
\end{exercise}

\begin{exercise}[$\star\star\star$]\label{exo:b2:nvs:8}
Sobre $E = C(\intcc{0}{1})$ con $\norm\cdot_\infty$, considera
\[
\varphi(f) = \sum_{n \geq 1} (-1)^n\, 2^{-n} f\bigl(\tfrac1n\bigr).
\]
Demuestra que $\varphi$ es una forma lineal \omterm{def:b2:metric:continuity}{continua} bien definida
con $\vertiii\varphi = 1$, pero que el supremo que define
$\vertiii\varphi$ \emph{no se alcanza} sobre la bola unidad cerrada.
\emph{(Cota superior: desigualdad triangular. \omterm{def:b2:nvs:norm}{Norma} $= 1$:
constrúyanse funciones \omterm{def:b2:metric:continuity}{continuas} $f_K$ con $\norm{f_K}_\infty \leq
1$ y $f_K(\frac1n) = (-1)^n$ para $n \leq K$; los puntos $\frac1n$
están aislados entre sí. No alcanzado: la igualdad obligaría a
$f(\frac1n) = (-1)^n$ para todo $n$, incompatible con la
\omterm{def:b2:metric:continuity}{continuidad} de $f$ en $0$, pues $\frac1n \to 0$.)}
\end{exercise}

\begin{exercise}[$\star\star\star$]\label{exo:b2:nvs:9}
Sea $E$ un espacio normado en el que la bola unidad cerrada es
\omterm{def:b2:metric:compact}{compacta}. Vuelve a deducir, sin citar el~\cref{thm:b2:nvs:riesz}, que toda sucesión acotada tiene una
subsucesión convergente, y demuestra que toda forma lineal sobre $E$
es \omterm{def:b2:metric:continuity}{continua} si y solo si $\dim E < \infty$. \emph{(Para la
dimensión infinita, constrúyase una forma discontinua definiéndola
libremente sobre una sucesión libre normalizada y extendiéndola,
admitiendo la existencia de un complementario algebraico.)}
\end{exercise}

\begin{exercise}[$\star\star$]\label{exo:b2:nvs:10}
Sobre $C(\intcc{0}{1})$, demuestra $\norm f_1 \leq \norm f_2 \leq
\norm f_\infty$ \emph{(Cauchy--Schwarz para la primera)} y prueba
con la familia $f_n(x) = x^n$ que ninguna de las dos desigualdades
puede invertirse salvo constante: las tres \omterm{def:b2:nvs:norm}{normas} son
no \omterm{def:b2:nvs:equivalent}{equivalentes} dos a dos.
\end{exercise}

\begin{exercise}[$\star\star$]\label{exo:b2:nvs:11}
(Distancia a un hiperplano) Sea $\varphi$ una forma lineal
\omterm{def:b2:metric:continuity}{continua} no nula sobre un espacio normado $E$. Demuestra que
\[
d\bigl(x, \ker\varphi\bigr) =
\frac{\abs{\varphi(x)}}{\vertiii\varphi}
\qquad (x \in E),
\]
y comprueba sobre el~\cref{exo:b2:nvs:8} que el ínfimo no tiene por
qué alcanzarse en ningún punto del hiperplano.
\end{exercise}

\begin{exercise}[$\star\star\star$]\label{exo:b2:nvs:12}
Sobre $E = \R[X]$ (todos los polinomios), sean $N_1(P) =
\sup_{\intcc{0}{1}}\abs P$ y $N_2(P) = \sup_{\intcc{0}{2}}\abs P$.
Prueba que $N_1 \leq N_2$ pero que $N_1$ y $N_2$ \emph{no} son
\omterm{def:b2:nvs:equivalent}{equivalentes}; deduce que la identidad $(E, N_2) \to (E, N_1)$ es
una biyección lineal \omterm{def:b2:metric:continuity}{continua} cuya inversa es discontinua. Prueba
por último que $(E, N_1)$ no es \omterm{def:b2:metric:complete}{completo} \emph{(sumas parciales de
Taylor de $\eu^x$)}. Los tres fenómenos son imposibles en dimensión
finita; di por qué.
\end{exercise}

\section{Problema: mejor aproximación y teorema de Chebyshev}

¿Con qué precisión puede aproximarse una función mediante polinomios
de grado dado, y qué polinomio lo hace mejor? Por el lado de la
existencia, la respuesta pertenece a este capítulo: la \omterm{def:b2:metric:compact}{compacidad} en
dimensión finita hace que las mejores aproximaciones existan. Por el
lado explícito, un caso no trivial se resuelve por completo a mano:
entre todos los polinomios \emph{mónicos} de grado $n$, el de menor
\omterm{def:b2:nvs:norm}{norma} del supremo sobre $\intcc{-1}{1}$ es el polinomio de Chebyshev
(normalizado), de \omterm{def:b2:nvs:norm}{norma} $2^{1-n}$: el \emph{teorema extremal de
Chebyshev}. El problema demuestra ambas caras y mide después con qué
gravedad falla la \omterm{def:b2:metric:compact}{compacidad} en dimensión infinita: la bola unidad de
$C(\intcc{0}{1})$ contiene constelaciones infinitas de puntos a
distancia mutua $1$.

\begin{problem}[{Problema de fin de semana --- el teorema extremal de
Chebyshev y la geometría de la bola unidad}]
\label{pb:b2:nvs:1}
Las \omterm{def:b2:nvs:norm}{normas} sin subíndice son \omterm{def:b2:nvs:norm}{normas} del supremo sobre el segmento
indicado.

\medskip
\noindent\textbf{Parte I --- Mejor aproximación en espacios
normados.}
\begin{enumerate}
  \item Sea $F$ un subespacio de dimensión finita de un espacio
        normado $E$ y $x \in E$. Demuestra que la distancia $d(x, F)
        = \inf_{f \in F}\norm{x - f}$ \emph{se alcanza}
        \emph{(redúcete a un subconjunto cerrado y acotado de $F$ y
        usa el~\cref{thm:b2:nvs:finitedim})}.
  \item Una \omterm{def:b2:nvs:norm}{norma} es \emph{estrictamente convexa} cuando $\norm u =
        \norm v = 1$ y $u \neq v$ implican $\bigl\Vert\frac{u +
        v}2\bigr\Vert < 1$. Prueba que $\norm\cdot_2$ sobre $\R^n$
        es estrictamente convexa \emph{(identidad del
        paralelogramo)} y que $\norm\cdot_1$ y $\norm\cdot_\infty$
        no lo son para $n \geq 2$.
  \item Demuestra que, para una \omterm{def:b2:nvs:norm}{norma} estrictamente convexa, la
        mejor aproximación de la pregunta 1 es \emph{única}.
  \item En $(\R^2, \norm\cdot_\infty)$, calcula todas las mejores
        aproximaciones de $x = (0, 1)$ por la recta $F =
        \operatorname{Vect}\bigl((1,0)\bigr)$: un intervalo de
        minimizadores.
  \item En $\bigl(C(\intcc{a}{b}), \norm\cdot_\infty\bigr)$, prueba
        que la mejor aproximación de $f$ por \emph{constantes} es
        única, igual a $c^* = \frac{\max f + \min f}{2}$, con
        distancia $\frac{\max f - \min f}{2}$; calcula ambas para
        $f(x) = x^2$ sobre $\intcc{0}{1}$.
\end{enumerate}

\medskip
\noindent\textbf{Parte II --- Polinomios de Chebyshev.}
\begin{enumerate}[resume]
  \item Prueba que existe exactamente un polinomio $T_n$ con
        $T_n(\cos\theta) = \cos n\theta$ para todo $\theta$
        \emph{(recurrencia $T_{n+1} = 2XT_n - T_{n-1}$ a partir de
        la fórmula de adición del coseno)}, que $\deg T_n = n$ y que
        su coeficiente director es $2^{n-1}$ para $n \geq 1$.
  \item Prueba que $\abs{T_n} \leq 1$ sobre $\intcc{-1}{1}$, con
        $T_n(\eta_k) = (-1)^k$ en los $n + 1$ puntos $\eta_k =
        \cos\frac{k\pi}{n}$ ($k = 0, \dots, n$), y que las raíces de
        $T_n$ son los $n$ puntos $\cos\frac{(2k-1)\pi}{2n}$,
        entrelazados con los $\eta_k$.
  \item Calcula $T_2, T_3, T_4$ y verifica la alternancia de $T_3$
        en $\eta_0, \dots, \eta_3 = 1, \frac12, -\frac12, -1$
        mediante evaluación directa.
  \item Para $\abs x \geq 1$, demuestra que
        \[
        T_n(x) = \frac{\bigl(x + \sqrt{x^2 - 1}\bigr)^n +
        \bigl(x - \sqrt{x^2 - 1}\bigr)^n}{2},
        \]
        y deduce que $T_n(x) \sim \frac12\bigl(x + \sqrt{x^2 -
        1}\bigr)^n \to \infty$ geométricamente para $x > 1$ fijo.
  \item Demuestra la ley de composición $T_m \circ T_n = T_{mn}$
        \emph{(compruébalo sobre $\intcc{-1}{1}$ e invoca la rigidez
        de los polinomios)}.
\end{enumerate}

\medskip
\noindent\textbf{Parte III --- El teorema extremal de Chebyshev.}
Escribamos $Q_n = 2^{1-n}T_n$ (mónico, por la pregunta 6).
\begin{enumerate}[resume]
  \item Sea $P$ mónico de grado $n \geq 1$ con
        $\sup_{\intcc{-1}{1}}\abs P < 2^{1-n}$. Evaluando $D = Q_n -
        P$ en los puntos $\eta_k$ y contando cambios de signo,
        obtén una contradicción. Concluye:
        \[
        \sup_{\intcc{-1}{1}}\abs P \;\geq\; 2^{1-n}
        \qquad\text{para todo } P \text{ mónico de grado } n.
        \]
  \item (Caso de igualdad) Supongamos $\sup_{\intcc{-1}{1}}\abs P =
        2^{1-n}$ con $P$ mónico de grado $n$, y sea $D = Q_n - P
        \neq 0$. Prueba que $(-1)^kD(\eta_k) \geq 0$ para todo $k$;
        prueba que cada uno de los $n$ intervalos
        $\intcc{\eta_{k}}{\eta_{k-1}}$ contiene un cero de $D$ y que
        un cero compartido por dos intervalos consecutivos es un
        punto \emph{interior} $\eta_k$ en el que además $D' = 0$.
        Concluye que $D$ tiene $n$ ceros contados con multiplicidad
        y, por tanto, $D = 0$: el minimizador es exactamente $Q_n$,
        el \emph{teorema extremal de Chebyshev}.
  \item Reformula el teorema como una distancia: sobre
        $\intcc{-1}{1}$,
        \[
        d_\infty\bigl(X^n,\ \R_{n-1}[X]\bigr) = 2^{1-n},
        \]
        con mejor aproximación única $X^n - Q_n$; y prueba mediante
        la sustitución afín $x = \frac{1+t}2$ que sobre
        $\intcc{0}{1}$ la distancia pasa a ser $2^{1-2n}$.
  \item (Nodos de interpolación óptimos) Para $n$ nodos $x_1, \dots,
        x_n \in \intcc{-1}{1}$, el polinomio nodal $\omega(x) =
        \prod_i(x - x_i)$ es mónico de grado $n$. Deduce de la
        pregunta 12 qué elección de nodos minimiza
        $\sup_{\intcc{-1}{1}}\abs\omega$, el factor dependiente de
        los nodos en la cota clásica del error de interpolación, y
        da el valor mínimo.
  \item Verifica a mano el caso $n = 2$ del teorema (halla
        directamente $\inf_c \sup_{\intcc{-1}{1}}\abs{x^2 - c}$) y
        calcula numéricamente la distancia de la pregunta 13 sobre
        $\intcc{0}{1}$ para $n = 10$. ¿Qué dice su tamaño sobre la
        gráfica de $x^{10}$?
\end{enumerate}

\medskip
\noindent\textbf{Parte IV --- La bola unidad de
$C(\intcc{0}{1})$.}
\begin{enumerate}[resume]
  \item Sea $g_k(x) = x^{2^k}$. Prueba que $\norm{g_k}_\infty = 1$ y
        que $\norm{g_k - g_j}_\infty \geq \frac14$ para $j > k$
        \emph{(evalúa en el punto donde $x^{2^k} = \frac12$)}: una
        sucesión acotada explícita sin subsucesión convergente; la
        bola unidad cerrada no es \omterm{def:b2:metric:compact}{compacta}, y a mano.
  \item (Lema de Riesz, afinado) Sea $F$ un \omterm{def:b2:reduction:eigen}{subespacio propio} de
        \emph{dimensión finita} de un espacio normado $E$. Usando la
        pregunta 1, produce un vector unitario $x$ con $d(x, F) = 1$
        exactamente, y no solo $\geq 1 - \varepsilon$ como en el
        lema del~\cref{thm:b2:nvs:riesz}.
  \item Deduce que en todo espacio normado de dimensión infinita hay
        una sucesión de vectores unitarios con distancias mutuas
        $\geq 1$, y vuelve a deducir de ahí el teorema de Riesz.
  \item En $C(\intcc{0}{1})$, exhibe explícitamente una
        constelación así: las funciones tienda $h_n$ con soporte en
        $\bigl[\frac1{n+1}, \frac1n\bigr]$ y valor máximo $1$.
        Verifica $\norm{h_n} = 1$ y $\norm{h_n - h_m} = 1$ para $n
        \neq m$, y observa que $h_n \to 0$ puntualmente pero no
        uniformemente.
  \item (Falla la precompacidad) Prueba que la bola unidad cerrada
        de $C(\intcc{0}{1})$ no puede recubrirse con un número
        finito de bolas de radio $\frac13$ \emph{(cada bola de ese
        tipo contiene a lo sumo un $h_n$)}; contrástese con el paso
        de precompacidad de la demostración del~\cref{thm:b2:metric:borellebesgue}.
\end{enumerate}

\medskip
\noindent\textbf{Parte V --- \omterm{def:b2:nvs:norm}{Normas} en acción sobre matrices, y
síntesis.}
\begin{enumerate}[resume]
  \item Demuestra que todo \omterm{def:b2:reduction:eigen}{valor propio} $\lambda$ de $A \in
        \mathcal{M}_n(\C)$ cumple $\abs\lambda \leq \vertiii A$ para
        toda \omterm{thm:b2:nvs:continuouslinear}{norma de operador}; aplica el~\cref{exo:b2:nvs:4} para
        acotar los \omterm{def:b2:reduction:eigen}{valores propios} de
        $\begin{psmallmatrix}1 & -2\\ 3 & 1\end{psmallmatrix}$ y
        compáralo con su módulo verdadero.
  \item (\omterm{def:b2:nvs:norm}{Normas} adaptadas) Sea $A$ \omterm{def:b2:reduction:diag}{diagonalizable}, $A =
        P\,\mathrm{diag}(\lambda_1, \dots, \lambda_n)\,P^{-1}$.
        Prueba que $N_P(x) = \norm{P^{-1}x}_\infty$ es una \omterm{def:b2:nvs:norm}{norma}
        cuya \omterm{thm:b2:nvs:continuouslinear}{norma de operador} cumple $\vertiii A_{N_P} =
        \max_i\abs{\lambda_i}$.
  \item Deduce que, para $A$ \omterm{def:b2:reduction:diag}{diagonalizable}, $A^k \to 0$ si y solo
        si todos los \omterm{def:b2:reduction:eigen}{valores propios} cumplen $\abs{\lambda_i} < 1$;
        la \omterm{thm:b2:nvs:finitedim}{equivalencia de normas} hace que la conclusión no dependa
        de la \omterm{def:b2:nvs:norm}{norma}. Compruébalo con $A =
        \frac14\begin{psmallmatrix}1 & 2\\ 2 &
        1\end{psmallmatrix}$.
  \item (Las constantes de equivalencia explotan) Sobre $\R_n[X]$,
        compara $N_c(P) = \max_k \abs{a_k}$ (coeficientes) con
        $\norm{P}_{\intcc{0}{1}}$: ambas son \omterm{def:b2:nvs:norm}{normas} y por tanto son
        \omterm{def:b2:nvs:equivalent}{equivalentes} para cada $n$ fijo; pero prueba, usando el
        minimizador mónico de la pregunta 13 sobre $\intcc{0}{1}$,
        que la mejor constante $C_n$ en $N_c \leq
        C_n\norm\cdot_{\intcc{0}{1}}$ cumple $C_n \geq 2^{2n-1}$.
        Concluye en una frase por qué ``todas las \omterm{def:b2:nvs:norm}{normas} son
        \omterm{def:b2:nvs:equivalent}{equivalentes}'' muere en dimensión infinita.
  \item (Síntesis) Una frase para cada punto: dónde trabajó la
        \omterm{def:b2:metric:compact}{compacidad} de las bolas de dimensión finita (preguntas 1 y
        12); qué gobierna la convexidad estricta; qué destruye la
        constelación de las preguntas 18--19; y cómo cuantifica la
        pregunta 24 ese fracaso. Nombra la cumbre (el teorema
        extremal de Chebyshev) e indica dónde encuentra la mejor
        aproximación su hogar moderno (el teorema de la proyección
        en espacios de Hilbert, volumen del tercer año, donde la
        \omterm{def:b2:metric:complete}{completitud} sustituye a la \omterm{def:b2:metric:compact}{compacidad}).
\end{enumerate}
\end{problem}
